% ============================================================================
% ESQUEMA: a sequência da metodologia como fluxo de decisões (HORIZONTAL v2).
% Encomenda do autor (tarefa 0138; reorientada na auditoria de 2026-08-24:
% "prefiro que fique na horizontal e menor, está ocupando uma página inteira").
% Conta a história da abertura do Cap. 3 na mesma ordem: problema -> terreno
% comum -> pilares 1-4 (setas com a lógica de precedência) -> LCE transversal
% -> fecho (validade e reprodutibilidade).
% Uso: abertura do Cap. 3, dentro de figure:
%   % ============================================================================
% ESQUEMA: a sequência da metodologia como fluxo de decisões (HORIZONTAL v2).
% Encomenda do autor (tarefa 0138; reorientada na auditoria de 2026-08-24:
% "prefiro que fique na horizontal e menor, está ocupando uma página inteira").
% Conta a história da abertura do Cap. 3 na mesma ordem: problema -> terreno
% comum -> pilares 1-4 (setas com a lógica de precedência) -> LCE transversal
% -> fecho (validade e reprodutibilidade).
% Uso: abertura do Cap. 3, dentro de figure:
%   % ============================================================================
% ESQUEMA: a sequência da metodologia como fluxo de decisões (HORIZONTAL v2).
% Encomenda do autor (tarefa 0138; reorientada na auditoria de 2026-08-24:
% "prefiro que fique na horizontal e menor, está ocupando uma página inteira").
% Conta a história da abertura do Cap. 3 na mesma ordem: problema -> terreno
% comum -> pilares 1-4 (setas com a lógica de precedência) -> LCE transversal
% -> fecho (validade e reprodutibilidade).
% Uso: abertura do Cap. 3, dentro de figure:
%   % ============================================================================
% ESQUEMA: a sequência da metodologia como fluxo de decisões (HORIZONTAL v2).
% Encomenda do autor (tarefa 0138; reorientada na auditoria de 2026-08-24:
% "prefiro que fique na horizontal e menor, está ocupando uma página inteira").
% Conta a história da abertura do Cap. 3 na mesma ordem: problema -> terreno
% comum -> pilares 1-4 (setas com a lógica de precedência) -> LCE transversal
% -> fecho (validade e reprodutibilidade).
% Uso: abertura do Cap. 3, dentro de figure:
%   \input{3-metodo/esquemas-propostos/esq-sequencia-metodologia}
% SEM números medidos, SEM códigos internos, SEM caminhos, SEM travessões.
% Uso standalone (pré-visualização):
%   \documentclass[tikz,border=6pt,12pt]{standalone}
%   \usetikzlibrary{arrows.meta, positioning}
%   \begin{document}\input{esq-sequencia-metodologia.tex}\end{document}
%   (no modo standalone, troque a \ref{ch:intro} por texto fixo)
% GEOMETRIA CALIBRADA PARA CORPO 12 (ppginf/book 12pt, textwidth 16cm):
% largura total ~15,6 cm, altura ~5 cm (fração de página, não página inteira).
% Legível em P&B: etapas = cinza claro; destino = cinza médio; LCE tracejada.
% ============================================================================
\begin{tikzpicture}[
    x=1cm, y=1cm,
    etapa/.style={draw, rounded corners=2pt, align=center, inner sep=2.5pt,
                  font=\scriptsize, fill=gray!8, minimum height=17mm},
    destino/.style={etapa, fill=gray!16},
    seta/.style={-{Stealth[length=2.2mm]}, semithick},
    rot/.style={font=\tiny, align=center, inner sep=1pt}
]
    % ------------------ a linha do processo, da esquerda para a direita ----
    \node[etapa, text width=20mm] (prob) at (0,0)
        {\textbf{problema}\\(Cap.~\ref{ch:intro}):\\reduzir o custo de\\rotulagem sem perda\\de desempenho};
    \node[etapa, text width=21mm] (terr) at (2.68,0)
        {\textbf{terreno}\\\textbf{comum}:\\desenho e registro;\\dados e auditoria;\\classificadores;\\métricas};
    \node[etapa, text width=20mm] (p1) at (5.36,0)
        {\textbf{pilar 1}:\\quanto importa a\\composição do\\conjunto inicial $L_0$};
    \node[etapa, text width=20mm] (p2) at (8.02,0)
        {\textbf{pilar 2}:\\construir $L_0$\\sem rótulos\\(DRI-SL)};
    \node[etapa, text width=20mm] (p3) at (10.66,0)
        {\textbf{pilar 3}:\\quem rotula:\\LLM como oráculo\\(acerto, custo, erro)};
    \node[destino, text width=21mm] (p4) at (13.42,0)
        {\textbf{pilar 4}:\\integração no\\FALCO e avaliação\\sob o mesmo\\orçamento};

    % ------------------ setas com a lógica de precedência ------------------
    \draw[seta] (prob) -- (terr);
    \draw[seta] (terr) -- (p1);
    \draw[seta] (p1) -- (p2);
    \draw[seta] (p2) -- (p3);
    \draw[seta] (p3) -- (p4);
    \node[rot, anchor=south] at (4.02,1.3) {o processo começa\\sem rótulo algum};
    \node[rot, anchor=south] at (6.69,1.3)  {medida a importância,\\constrói-se o começo};
    \node[rot, anchor=south] at (9.34,1.3)  {definido o começo:\\quem rotula?};
    \node[rot, anchor=south] at (12.04,1.3)  {resolvidos início\\e oráculo};

    % ------------------ LCE transversal + fecho, abaixo da linha -----------
    \draw[dashed, rounded corners=2pt]
        (4.2,-1.55) rectangle (14.55,-2.35);
    \node[font=\scriptsize, align=center] at (9.37,-1.95)
        {\textbf{transversal}: a LCE resume as curvas de aprendizado dos quatro pilares};
    \node[font=\tiny, align=center, anchor=north] at (7.25,-2.6)
        {o capítulo encerra com a síntese do desenho experimental};
\end{tikzpicture}

% SEM números medidos, SEM códigos internos, SEM caminhos, SEM travessões.
% Uso standalone (pré-visualização):
%   \documentclass[tikz,border=6pt,12pt]{standalone}
%   \usetikzlibrary{arrows.meta, positioning}
%   \begin{document}% ============================================================================
% ESQUEMA: a sequência da metodologia como fluxo de decisões (HORIZONTAL v2).
% Encomenda do autor (tarefa 0138; reorientada na auditoria de 2026-08-24:
% "prefiro que fique na horizontal e menor, está ocupando uma página inteira").
% Conta a história da abertura do Cap. 3 na mesma ordem: problema -> terreno
% comum -> pilares 1-4 (setas com a lógica de precedência) -> LCE transversal
% -> fecho (validade e reprodutibilidade).
% Uso: abertura do Cap. 3, dentro de figure:
%   \input{3-metodo/esquemas-propostos/esq-sequencia-metodologia}
% SEM números medidos, SEM códigos internos, SEM caminhos, SEM travessões.
% Uso standalone (pré-visualização):
%   \documentclass[tikz,border=6pt,12pt]{standalone}
%   \usetikzlibrary{arrows.meta, positioning}
%   \begin{document}\input{esq-sequencia-metodologia.tex}\end{document}
%   (no modo standalone, troque a \ref{ch:intro} por texto fixo)
% GEOMETRIA CALIBRADA PARA CORPO 12 (ppginf/book 12pt, textwidth 16cm):
% largura total ~15,6 cm, altura ~5 cm (fração de página, não página inteira).
% Legível em P&B: etapas = cinza claro; destino = cinza médio; LCE tracejada.
% ============================================================================
\begin{tikzpicture}[
    x=1cm, y=1cm,
    etapa/.style={draw, rounded corners=2pt, align=center, inner sep=2.5pt,
                  font=\scriptsize, fill=gray!8, minimum height=17mm},
    destino/.style={etapa, fill=gray!16},
    seta/.style={-{Stealth[length=2.2mm]}, semithick},
    rot/.style={font=\tiny, align=center, inner sep=1pt}
]
    % ------------------ a linha do processo, da esquerda para a direita ----
    \node[etapa, text width=20mm] (prob) at (0,0)
        {\textbf{problema}\\(Cap.~\ref{ch:intro}):\\reduzir o custo de\\rotulagem sem perda\\de desempenho};
    \node[etapa, text width=21mm] (terr) at (2.68,0)
        {\textbf{terreno}\\\textbf{comum}:\\desenho e registro;\\dados e auditoria;\\classificadores;\\métricas};
    \node[etapa, text width=20mm] (p1) at (5.36,0)
        {\textbf{pilar 1}:\\quanto importa a\\composição do\\conjunto inicial $L_0$};
    \node[etapa, text width=20mm] (p2) at (8.02,0)
        {\textbf{pilar 2}:\\construir $L_0$\\sem rótulos\\(DRI-SL)};
    \node[etapa, text width=20mm] (p3) at (10.66,0)
        {\textbf{pilar 3}:\\quem rotula:\\LLM como oráculo\\(acerto, custo, erro)};
    \node[destino, text width=21mm] (p4) at (13.42,0)
        {\textbf{pilar 4}:\\integração no\\FALCO e avaliação\\sob o mesmo\\orçamento};

    % ------------------ setas com a lógica de precedência ------------------
    \draw[seta] (prob) -- (terr);
    \draw[seta] (terr) -- (p1);
    \draw[seta] (p1) -- (p2);
    \draw[seta] (p2) -- (p3);
    \draw[seta] (p3) -- (p4);
    \node[rot, anchor=south] at (4.02,1.3) {o processo começa\\sem rótulo algum};
    \node[rot, anchor=south] at (6.69,1.3)  {medida a importância,\\constrói-se o começo};
    \node[rot, anchor=south] at (9.34,1.3)  {definido o começo:\\quem rotula?};
    \node[rot, anchor=south] at (12.04,1.3)  {resolvidos início\\e oráculo};

    % ------------------ LCE transversal + fecho, abaixo da linha -----------
    \draw[dashed, rounded corners=2pt]
        (4.2,-1.55) rectangle (14.55,-2.35);
    \node[font=\scriptsize, align=center] at (9.37,-1.95)
        {\textbf{transversal}: a LCE resume as curvas de aprendizado dos quatro pilares};
    \node[font=\tiny, align=center, anchor=north] at (7.25,-2.6)
        {o capítulo encerra com a síntese do desenho experimental};
\end{tikzpicture}
\end{document}
%   (no modo standalone, troque a \ref{ch:intro} por texto fixo)
% GEOMETRIA CALIBRADA PARA CORPO 12 (ppginf/book 12pt, textwidth 16cm):
% largura total ~15,6 cm, altura ~5 cm (fração de página, não página inteira).
% Legível em P&B: etapas = cinza claro; destino = cinza médio; LCE tracejada.
% ============================================================================
\begin{tikzpicture}[
    x=1cm, y=1cm,
    etapa/.style={draw, rounded corners=2pt, align=center, inner sep=2.5pt,
                  font=\scriptsize, fill=gray!8, minimum height=17mm},
    destino/.style={etapa, fill=gray!16},
    seta/.style={-{Stealth[length=2.2mm]}, semithick},
    rot/.style={font=\tiny, align=center, inner sep=1pt}
]
    % ------------------ a linha do processo, da esquerda para a direita ----
    \node[etapa, text width=20mm] (prob) at (0,0)
        {\textbf{problema}\\(Cap.~\ref{ch:intro}):\\reduzir o custo de\\rotulagem sem perda\\de desempenho};
    \node[etapa, text width=21mm] (terr) at (2.68,0)
        {\textbf{terreno}\\\textbf{comum}:\\desenho e registro;\\dados e auditoria;\\classificadores;\\métricas};
    \node[etapa, text width=20mm] (p1) at (5.36,0)
        {\textbf{pilar 1}:\\quanto importa a\\composição do\\conjunto inicial $L_0$};
    \node[etapa, text width=20mm] (p2) at (8.02,0)
        {\textbf{pilar 2}:\\construir $L_0$\\sem rótulos\\(DRI-SL)};
    \node[etapa, text width=20mm] (p3) at (10.66,0)
        {\textbf{pilar 3}:\\quem rotula:\\LLM como oráculo\\(acerto, custo, erro)};
    \node[destino, text width=21mm] (p4) at (13.42,0)
        {\textbf{pilar 4}:\\integração no\\FALCO e avaliação\\sob o mesmo\\orçamento};

    % ------------------ setas com a lógica de precedência ------------------
    \draw[seta] (prob) -- (terr);
    \draw[seta] (terr) -- (p1);
    \draw[seta] (p1) -- (p2);
    \draw[seta] (p2) -- (p3);
    \draw[seta] (p3) -- (p4);
    \node[rot, anchor=south] at (4.02,1.3) {o processo começa\\sem rótulo algum};
    \node[rot, anchor=south] at (6.69,1.3)  {medida a importância,\\constrói-se o começo};
    \node[rot, anchor=south] at (9.34,1.3)  {definido o começo:\\quem rotula?};
    \node[rot, anchor=south] at (12.04,1.3)  {resolvidos início\\e oráculo};

    % ------------------ LCE transversal + fecho, abaixo da linha -----------
    \draw[dashed, rounded corners=2pt]
        (4.2,-1.55) rectangle (14.55,-2.35);
    \node[font=\scriptsize, align=center] at (9.37,-1.95)
        {\textbf{transversal}: a LCE resume as curvas de aprendizado dos quatro pilares};
    \node[font=\tiny, align=center, anchor=north] at (7.25,-2.6)
        {o capítulo encerra com a síntese do desenho experimental};
\end{tikzpicture}

% SEM números medidos, SEM códigos internos, SEM caminhos, SEM travessões.
% Uso standalone (pré-visualização):
%   \documentclass[tikz,border=6pt,12pt]{standalone}
%   \usetikzlibrary{arrows.meta, positioning}
%   \begin{document}% ============================================================================
% ESQUEMA: a sequência da metodologia como fluxo de decisões (HORIZONTAL v2).
% Encomenda do autor (tarefa 0138; reorientada na auditoria de 2026-08-24:
% "prefiro que fique na horizontal e menor, está ocupando uma página inteira").
% Conta a história da abertura do Cap. 3 na mesma ordem: problema -> terreno
% comum -> pilares 1-4 (setas com a lógica de precedência) -> LCE transversal
% -> fecho (validade e reprodutibilidade).
% Uso: abertura do Cap. 3, dentro de figure:
%   % ============================================================================
% ESQUEMA: a sequência da metodologia como fluxo de decisões (HORIZONTAL v2).
% Encomenda do autor (tarefa 0138; reorientada na auditoria de 2026-08-24:
% "prefiro que fique na horizontal e menor, está ocupando uma página inteira").
% Conta a história da abertura do Cap. 3 na mesma ordem: problema -> terreno
% comum -> pilares 1-4 (setas com a lógica de precedência) -> LCE transversal
% -> fecho (validade e reprodutibilidade).
% Uso: abertura do Cap. 3, dentro de figure:
%   \input{3-metodo/esquemas-propostos/esq-sequencia-metodologia}
% SEM números medidos, SEM códigos internos, SEM caminhos, SEM travessões.
% Uso standalone (pré-visualização):
%   \documentclass[tikz,border=6pt,12pt]{standalone}
%   \usetikzlibrary{arrows.meta, positioning}
%   \begin{document}\input{esq-sequencia-metodologia.tex}\end{document}
%   (no modo standalone, troque a \ref{ch:intro} por texto fixo)
% GEOMETRIA CALIBRADA PARA CORPO 12 (ppginf/book 12pt, textwidth 16cm):
% largura total ~15,6 cm, altura ~5 cm (fração de página, não página inteira).
% Legível em P&B: etapas = cinza claro; destino = cinza médio; LCE tracejada.
% ============================================================================
\begin{tikzpicture}[
    x=1cm, y=1cm,
    etapa/.style={draw, rounded corners=2pt, align=center, inner sep=2.5pt,
                  font=\scriptsize, fill=gray!8, minimum height=17mm},
    destino/.style={etapa, fill=gray!16},
    seta/.style={-{Stealth[length=2.2mm]}, semithick},
    rot/.style={font=\tiny, align=center, inner sep=1pt}
]
    % ------------------ a linha do processo, da esquerda para a direita ----
    \node[etapa, text width=20mm] (prob) at (0,0)
        {\textbf{problema}\\(Cap.~\ref{ch:intro}):\\reduzir o custo de\\rotulagem sem perda\\de desempenho};
    \node[etapa, text width=21mm] (terr) at (2.68,0)
        {\textbf{terreno}\\\textbf{comum}:\\desenho e registro;\\dados e auditoria;\\classificadores;\\métricas};
    \node[etapa, text width=20mm] (p1) at (5.36,0)
        {\textbf{pilar 1}:\\quanto importa a\\composição do\\conjunto inicial $L_0$};
    \node[etapa, text width=20mm] (p2) at (8.02,0)
        {\textbf{pilar 2}:\\construir $L_0$\\sem rótulos\\(DRI-SL)};
    \node[etapa, text width=20mm] (p3) at (10.66,0)
        {\textbf{pilar 3}:\\quem rotula:\\LLM como oráculo\\(acerto, custo, erro)};
    \node[destino, text width=21mm] (p4) at (13.42,0)
        {\textbf{pilar 4}:\\integração no\\FALCO e avaliação\\sob o mesmo\\orçamento};

    % ------------------ setas com a lógica de precedência ------------------
    \draw[seta] (prob) -- (terr);
    \draw[seta] (terr) -- (p1);
    \draw[seta] (p1) -- (p2);
    \draw[seta] (p2) -- (p3);
    \draw[seta] (p3) -- (p4);
    \node[rot, anchor=south] at (4.02,1.3) {o processo começa\\sem rótulo algum};
    \node[rot, anchor=south] at (6.69,1.3)  {medida a importância,\\constrói-se o começo};
    \node[rot, anchor=south] at (9.34,1.3)  {definido o começo:\\quem rotula?};
    \node[rot, anchor=south] at (12.04,1.3)  {resolvidos início\\e oráculo};

    % ------------------ LCE transversal + fecho, abaixo da linha -----------
    \draw[dashed, rounded corners=2pt]
        (4.2,-1.55) rectangle (14.55,-2.35);
    \node[font=\scriptsize, align=center] at (9.37,-1.95)
        {\textbf{transversal}: a LCE resume as curvas de aprendizado dos quatro pilares};
    \node[font=\tiny, align=center, anchor=north] at (7.25,-2.6)
        {o capítulo encerra com a síntese do desenho experimental};
\end{tikzpicture}

% SEM números medidos, SEM códigos internos, SEM caminhos, SEM travessões.
% Uso standalone (pré-visualização):
%   \documentclass[tikz,border=6pt,12pt]{standalone}
%   \usetikzlibrary{arrows.meta, positioning}
%   \begin{document}% ============================================================================
% ESQUEMA: a sequência da metodologia como fluxo de decisões (HORIZONTAL v2).
% Encomenda do autor (tarefa 0138; reorientada na auditoria de 2026-08-24:
% "prefiro que fique na horizontal e menor, está ocupando uma página inteira").
% Conta a história da abertura do Cap. 3 na mesma ordem: problema -> terreno
% comum -> pilares 1-4 (setas com a lógica de precedência) -> LCE transversal
% -> fecho (validade e reprodutibilidade).
% Uso: abertura do Cap. 3, dentro de figure:
%   \input{3-metodo/esquemas-propostos/esq-sequencia-metodologia}
% SEM números medidos, SEM códigos internos, SEM caminhos, SEM travessões.
% Uso standalone (pré-visualização):
%   \documentclass[tikz,border=6pt,12pt]{standalone}
%   \usetikzlibrary{arrows.meta, positioning}
%   \begin{document}\input{esq-sequencia-metodologia.tex}\end{document}
%   (no modo standalone, troque a \ref{ch:intro} por texto fixo)
% GEOMETRIA CALIBRADA PARA CORPO 12 (ppginf/book 12pt, textwidth 16cm):
% largura total ~15,6 cm, altura ~5 cm (fração de página, não página inteira).
% Legível em P&B: etapas = cinza claro; destino = cinza médio; LCE tracejada.
% ============================================================================
\begin{tikzpicture}[
    x=1cm, y=1cm,
    etapa/.style={draw, rounded corners=2pt, align=center, inner sep=2.5pt,
                  font=\scriptsize, fill=gray!8, minimum height=17mm},
    destino/.style={etapa, fill=gray!16},
    seta/.style={-{Stealth[length=2.2mm]}, semithick},
    rot/.style={font=\tiny, align=center, inner sep=1pt}
]
    % ------------------ a linha do processo, da esquerda para a direita ----
    \node[etapa, text width=20mm] (prob) at (0,0)
        {\textbf{problema}\\(Cap.~\ref{ch:intro}):\\reduzir o custo de\\rotulagem sem perda\\de desempenho};
    \node[etapa, text width=21mm] (terr) at (2.68,0)
        {\textbf{terreno}\\\textbf{comum}:\\desenho e registro;\\dados e auditoria;\\classificadores;\\métricas};
    \node[etapa, text width=20mm] (p1) at (5.36,0)
        {\textbf{pilar 1}:\\quanto importa a\\composição do\\conjunto inicial $L_0$};
    \node[etapa, text width=20mm] (p2) at (8.02,0)
        {\textbf{pilar 2}:\\construir $L_0$\\sem rótulos\\(DRI-SL)};
    \node[etapa, text width=20mm] (p3) at (10.66,0)
        {\textbf{pilar 3}:\\quem rotula:\\LLM como oráculo\\(acerto, custo, erro)};
    \node[destino, text width=21mm] (p4) at (13.42,0)
        {\textbf{pilar 4}:\\integração no\\FALCO e avaliação\\sob o mesmo\\orçamento};

    % ------------------ setas com a lógica de precedência ------------------
    \draw[seta] (prob) -- (terr);
    \draw[seta] (terr) -- (p1);
    \draw[seta] (p1) -- (p2);
    \draw[seta] (p2) -- (p3);
    \draw[seta] (p3) -- (p4);
    \node[rot, anchor=south] at (4.02,1.3) {o processo começa\\sem rótulo algum};
    \node[rot, anchor=south] at (6.69,1.3)  {medida a importância,\\constrói-se o começo};
    \node[rot, anchor=south] at (9.34,1.3)  {definido o começo:\\quem rotula?};
    \node[rot, anchor=south] at (12.04,1.3)  {resolvidos início\\e oráculo};

    % ------------------ LCE transversal + fecho, abaixo da linha -----------
    \draw[dashed, rounded corners=2pt]
        (4.2,-1.55) rectangle (14.55,-2.35);
    \node[font=\scriptsize, align=center] at (9.37,-1.95)
        {\textbf{transversal}: a LCE resume as curvas de aprendizado dos quatro pilares};
    \node[font=\tiny, align=center, anchor=north] at (7.25,-2.6)
        {o capítulo encerra com a síntese do desenho experimental};
\end{tikzpicture}
\end{document}
%   (no modo standalone, troque a \ref{ch:intro} por texto fixo)
% GEOMETRIA CALIBRADA PARA CORPO 12 (ppginf/book 12pt, textwidth 16cm):
% largura total ~15,6 cm, altura ~5 cm (fração de página, não página inteira).
% Legível em P&B: etapas = cinza claro; destino = cinza médio; LCE tracejada.
% ============================================================================
\begin{tikzpicture}[
    x=1cm, y=1cm,
    etapa/.style={draw, rounded corners=2pt, align=center, inner sep=2.5pt,
                  font=\scriptsize, fill=gray!8, minimum height=17mm},
    destino/.style={etapa, fill=gray!16},
    seta/.style={-{Stealth[length=2.2mm]}, semithick},
    rot/.style={font=\tiny, align=center, inner sep=1pt}
]
    % ------------------ a linha do processo, da esquerda para a direita ----
    \node[etapa, text width=20mm] (prob) at (0,0)
        {\textbf{problema}\\(Cap.~\ref{ch:intro}):\\reduzir o custo de\\rotulagem sem perda\\de desempenho};
    \node[etapa, text width=21mm] (terr) at (2.68,0)
        {\textbf{terreno}\\\textbf{comum}:\\desenho e registro;\\dados e auditoria;\\classificadores;\\métricas};
    \node[etapa, text width=20mm] (p1) at (5.36,0)
        {\textbf{pilar 1}:\\quanto importa a\\composição do\\conjunto inicial $L_0$};
    \node[etapa, text width=20mm] (p2) at (8.02,0)
        {\textbf{pilar 2}:\\construir $L_0$\\sem rótulos\\(DRI-SL)};
    \node[etapa, text width=20mm] (p3) at (10.66,0)
        {\textbf{pilar 3}:\\quem rotula:\\LLM como oráculo\\(acerto, custo, erro)};
    \node[destino, text width=21mm] (p4) at (13.42,0)
        {\textbf{pilar 4}:\\integração no\\FALCO e avaliação\\sob o mesmo\\orçamento};

    % ------------------ setas com a lógica de precedência ------------------
    \draw[seta] (prob) -- (terr);
    \draw[seta] (terr) -- (p1);
    \draw[seta] (p1) -- (p2);
    \draw[seta] (p2) -- (p3);
    \draw[seta] (p3) -- (p4);
    \node[rot, anchor=south] at (4.02,1.3) {o processo começa\\sem rótulo algum};
    \node[rot, anchor=south] at (6.69,1.3)  {medida a importância,\\constrói-se o começo};
    \node[rot, anchor=south] at (9.34,1.3)  {definido o começo:\\quem rotula?};
    \node[rot, anchor=south] at (12.04,1.3)  {resolvidos início\\e oráculo};

    % ------------------ LCE transversal + fecho, abaixo da linha -----------
    \draw[dashed, rounded corners=2pt]
        (4.2,-1.55) rectangle (14.55,-2.35);
    \node[font=\scriptsize, align=center] at (9.37,-1.95)
        {\textbf{transversal}: a LCE resume as curvas de aprendizado dos quatro pilares};
    \node[font=\tiny, align=center, anchor=north] at (7.25,-2.6)
        {o capítulo encerra com a síntese do desenho experimental};
\end{tikzpicture}
\end{document}
%   (no modo standalone, troque a \ref{ch:intro} por texto fixo)
% GEOMETRIA CALIBRADA PARA CORPO 12 (ppginf/book 12pt, textwidth 16cm):
% largura total ~15,6 cm, altura ~5 cm (fração de página, não página inteira).
% Legível em P&B: etapas = cinza claro; destino = cinza médio; LCE tracejada.
% ============================================================================
\begin{tikzpicture}[
    x=1cm, y=1cm,
    etapa/.style={draw, rounded corners=2pt, align=center, inner sep=2.5pt,
                  font=\scriptsize, fill=gray!8, minimum height=17mm},
    destino/.style={etapa, fill=gray!16},
    seta/.style={-{Stealth[length=2.2mm]}, semithick},
    rot/.style={font=\tiny, align=center, inner sep=1pt}
]
    % ------------------ a linha do processo, da esquerda para a direita ----
    \node[etapa, text width=20mm] (prob) at (0,0)
        {\textbf{problema}\\(Cap.~\ref{ch:intro}):\\reduzir o custo de\\rotulagem sem perda\\de desempenho};
    \node[etapa, text width=21mm] (terr) at (2.68,0)
        {\textbf{terreno}\\\textbf{comum}:\\desenho e registro;\\dados e auditoria;\\classificadores;\\métricas};
    \node[etapa, text width=20mm] (p1) at (5.36,0)
        {\textbf{pilar 1}:\\quanto importa a\\composição do\\conjunto inicial $L_0$};
    \node[etapa, text width=20mm] (p2) at (8.02,0)
        {\textbf{pilar 2}:\\construir $L_0$\\sem rótulos\\(DRI-SL)};
    \node[etapa, text width=20mm] (p3) at (10.66,0)
        {\textbf{pilar 3}:\\quem rotula:\\LLM como oráculo\\(acerto, custo, erro)};
    \node[destino, text width=21mm] (p4) at (13.42,0)
        {\textbf{pilar 4}:\\integração no\\FALCO e avaliação\\sob o mesmo\\orçamento};

    % ------------------ setas com a lógica de precedência ------------------
    \draw[seta] (prob) -- (terr);
    \draw[seta] (terr) -- (p1);
    \draw[seta] (p1) -- (p2);
    \draw[seta] (p2) -- (p3);
    \draw[seta] (p3) -- (p4);
    \node[rot, anchor=south] at (4.02,1.3) {o processo começa\\sem rótulo algum};
    \node[rot, anchor=south] at (6.69,1.3)  {medida a importância,\\constrói-se o começo};
    \node[rot, anchor=south] at (9.34,1.3)  {definido o começo:\\quem rotula?};
    \node[rot, anchor=south] at (12.04,1.3)  {resolvidos início\\e oráculo};

    % ------------------ LCE transversal + fecho, abaixo da linha -----------
    \draw[dashed, rounded corners=2pt]
        (4.2,-1.55) rectangle (14.55,-2.35);
    \node[font=\scriptsize, align=center] at (9.37,-1.95)
        {\textbf{transversal}: a LCE resume as curvas de aprendizado dos quatro pilares};
    \node[font=\tiny, align=center, anchor=north] at (7.25,-2.6)
        {o capítulo encerra com a síntese do desenho experimental};
\end{tikzpicture}

% SEM números medidos, SEM códigos internos, SEM caminhos, SEM travessões.
% Uso standalone (pré-visualização):
%   \documentclass[tikz,border=6pt,12pt]{standalone}
%   \usetikzlibrary{arrows.meta, positioning}
%   \begin{document}% ============================================================================
% ESQUEMA: a sequência da metodologia como fluxo de decisões (HORIZONTAL v2).
% Encomenda do autor (tarefa 0138; reorientada na auditoria de 2026-08-24:
% "prefiro que fique na horizontal e menor, está ocupando uma página inteira").
% Conta a história da abertura do Cap. 3 na mesma ordem: problema -> terreno
% comum -> pilares 1-4 (setas com a lógica de precedência) -> LCE transversal
% -> fecho (validade e reprodutibilidade).
% Uso: abertura do Cap. 3, dentro de figure:
%   % ============================================================================
% ESQUEMA: a sequência da metodologia como fluxo de decisões (HORIZONTAL v2).
% Encomenda do autor (tarefa 0138; reorientada na auditoria de 2026-08-24:
% "prefiro que fique na horizontal e menor, está ocupando uma página inteira").
% Conta a história da abertura do Cap. 3 na mesma ordem: problema -> terreno
% comum -> pilares 1-4 (setas com a lógica de precedência) -> LCE transversal
% -> fecho (validade e reprodutibilidade).
% Uso: abertura do Cap. 3, dentro de figure:
%   % ============================================================================
% ESQUEMA: a sequência da metodologia como fluxo de decisões (HORIZONTAL v2).
% Encomenda do autor (tarefa 0138; reorientada na auditoria de 2026-08-24:
% "prefiro que fique na horizontal e menor, está ocupando uma página inteira").
% Conta a história da abertura do Cap. 3 na mesma ordem: problema -> terreno
% comum -> pilares 1-4 (setas com a lógica de precedência) -> LCE transversal
% -> fecho (validade e reprodutibilidade).
% Uso: abertura do Cap. 3, dentro de figure:
%   \input{3-metodo/esquemas-propostos/esq-sequencia-metodologia}
% SEM números medidos, SEM códigos internos, SEM caminhos, SEM travessões.
% Uso standalone (pré-visualização):
%   \documentclass[tikz,border=6pt,12pt]{standalone}
%   \usetikzlibrary{arrows.meta, positioning}
%   \begin{document}\input{esq-sequencia-metodologia.tex}\end{document}
%   (no modo standalone, troque a \ref{ch:intro} por texto fixo)
% GEOMETRIA CALIBRADA PARA CORPO 12 (ppginf/book 12pt, textwidth 16cm):
% largura total ~15,6 cm, altura ~5 cm (fração de página, não página inteira).
% Legível em P&B: etapas = cinza claro; destino = cinza médio; LCE tracejada.
% ============================================================================
\begin{tikzpicture}[
    x=1cm, y=1cm,
    etapa/.style={draw, rounded corners=2pt, align=center, inner sep=2.5pt,
                  font=\scriptsize, fill=gray!8, minimum height=17mm},
    destino/.style={etapa, fill=gray!16},
    seta/.style={-{Stealth[length=2.2mm]}, semithick},
    rot/.style={font=\tiny, align=center, inner sep=1pt}
]
    % ------------------ a linha do processo, da esquerda para a direita ----
    \node[etapa, text width=20mm] (prob) at (0,0)
        {\textbf{problema}\\(Cap.~\ref{ch:intro}):\\reduzir o custo de\\rotulagem sem perda\\de desempenho};
    \node[etapa, text width=21mm] (terr) at (2.68,0)
        {\textbf{terreno}\\\textbf{comum}:\\desenho e registro;\\dados e auditoria;\\classificadores;\\métricas};
    \node[etapa, text width=20mm] (p1) at (5.36,0)
        {\textbf{pilar 1}:\\quanto importa a\\composição do\\conjunto inicial $L_0$};
    \node[etapa, text width=20mm] (p2) at (8.02,0)
        {\textbf{pilar 2}:\\construir $L_0$\\sem rótulos\\(DRI-SL)};
    \node[etapa, text width=20mm] (p3) at (10.66,0)
        {\textbf{pilar 3}:\\quem rotula:\\LLM como oráculo\\(acerto, custo, erro)};
    \node[destino, text width=21mm] (p4) at (13.42,0)
        {\textbf{pilar 4}:\\integração no\\FALCO e avaliação\\sob o mesmo\\orçamento};

    % ------------------ setas com a lógica de precedência ------------------
    \draw[seta] (prob) -- (terr);
    \draw[seta] (terr) -- (p1);
    \draw[seta] (p1) -- (p2);
    \draw[seta] (p2) -- (p3);
    \draw[seta] (p3) -- (p4);
    \node[rot, anchor=south] at (4.02,1.3) {o processo começa\\sem rótulo algum};
    \node[rot, anchor=south] at (6.69,1.3)  {medida a importância,\\constrói-se o começo};
    \node[rot, anchor=south] at (9.34,1.3)  {definido o começo:\\quem rotula?};
    \node[rot, anchor=south] at (12.04,1.3)  {resolvidos início\\e oráculo};

    % ------------------ LCE transversal + fecho, abaixo da linha -----------
    \draw[dashed, rounded corners=2pt]
        (4.2,-1.55) rectangle (14.55,-2.35);
    \node[font=\scriptsize, align=center] at (9.37,-1.95)
        {\textbf{transversal}: a LCE resume as curvas de aprendizado dos quatro pilares};
    \node[font=\tiny, align=center, anchor=north] at (7.25,-2.6)
        {o capítulo encerra com a síntese do desenho experimental};
\end{tikzpicture}

% SEM números medidos, SEM códigos internos, SEM caminhos, SEM travessões.
% Uso standalone (pré-visualização):
%   \documentclass[tikz,border=6pt,12pt]{standalone}
%   \usetikzlibrary{arrows.meta, positioning}
%   \begin{document}% ============================================================================
% ESQUEMA: a sequência da metodologia como fluxo de decisões (HORIZONTAL v2).
% Encomenda do autor (tarefa 0138; reorientada na auditoria de 2026-08-24:
% "prefiro que fique na horizontal e menor, está ocupando uma página inteira").
% Conta a história da abertura do Cap. 3 na mesma ordem: problema -> terreno
% comum -> pilares 1-4 (setas com a lógica de precedência) -> LCE transversal
% -> fecho (validade e reprodutibilidade).
% Uso: abertura do Cap. 3, dentro de figure:
%   \input{3-metodo/esquemas-propostos/esq-sequencia-metodologia}
% SEM números medidos, SEM códigos internos, SEM caminhos, SEM travessões.
% Uso standalone (pré-visualização):
%   \documentclass[tikz,border=6pt,12pt]{standalone}
%   \usetikzlibrary{arrows.meta, positioning}
%   \begin{document}\input{esq-sequencia-metodologia.tex}\end{document}
%   (no modo standalone, troque a \ref{ch:intro} por texto fixo)
% GEOMETRIA CALIBRADA PARA CORPO 12 (ppginf/book 12pt, textwidth 16cm):
% largura total ~15,6 cm, altura ~5 cm (fração de página, não página inteira).
% Legível em P&B: etapas = cinza claro; destino = cinza médio; LCE tracejada.
% ============================================================================
\begin{tikzpicture}[
    x=1cm, y=1cm,
    etapa/.style={draw, rounded corners=2pt, align=center, inner sep=2.5pt,
                  font=\scriptsize, fill=gray!8, minimum height=17mm},
    destino/.style={etapa, fill=gray!16},
    seta/.style={-{Stealth[length=2.2mm]}, semithick},
    rot/.style={font=\tiny, align=center, inner sep=1pt}
]
    % ------------------ a linha do processo, da esquerda para a direita ----
    \node[etapa, text width=20mm] (prob) at (0,0)
        {\textbf{problema}\\(Cap.~\ref{ch:intro}):\\reduzir o custo de\\rotulagem sem perda\\de desempenho};
    \node[etapa, text width=21mm] (terr) at (2.68,0)
        {\textbf{terreno}\\\textbf{comum}:\\desenho e registro;\\dados e auditoria;\\classificadores;\\métricas};
    \node[etapa, text width=20mm] (p1) at (5.36,0)
        {\textbf{pilar 1}:\\quanto importa a\\composição do\\conjunto inicial $L_0$};
    \node[etapa, text width=20mm] (p2) at (8.02,0)
        {\textbf{pilar 2}:\\construir $L_0$\\sem rótulos\\(DRI-SL)};
    \node[etapa, text width=20mm] (p3) at (10.66,0)
        {\textbf{pilar 3}:\\quem rotula:\\LLM como oráculo\\(acerto, custo, erro)};
    \node[destino, text width=21mm] (p4) at (13.42,0)
        {\textbf{pilar 4}:\\integração no\\FALCO e avaliação\\sob o mesmo\\orçamento};

    % ------------------ setas com a lógica de precedência ------------------
    \draw[seta] (prob) -- (terr);
    \draw[seta] (terr) -- (p1);
    \draw[seta] (p1) -- (p2);
    \draw[seta] (p2) -- (p3);
    \draw[seta] (p3) -- (p4);
    \node[rot, anchor=south] at (4.02,1.3) {o processo começa\\sem rótulo algum};
    \node[rot, anchor=south] at (6.69,1.3)  {medida a importância,\\constrói-se o começo};
    \node[rot, anchor=south] at (9.34,1.3)  {definido o começo:\\quem rotula?};
    \node[rot, anchor=south] at (12.04,1.3)  {resolvidos início\\e oráculo};

    % ------------------ LCE transversal + fecho, abaixo da linha -----------
    \draw[dashed, rounded corners=2pt]
        (4.2,-1.55) rectangle (14.55,-2.35);
    \node[font=\scriptsize, align=center] at (9.37,-1.95)
        {\textbf{transversal}: a LCE resume as curvas de aprendizado dos quatro pilares};
    \node[font=\tiny, align=center, anchor=north] at (7.25,-2.6)
        {o capítulo encerra com a síntese do desenho experimental};
\end{tikzpicture}
\end{document}
%   (no modo standalone, troque a \ref{ch:intro} por texto fixo)
% GEOMETRIA CALIBRADA PARA CORPO 12 (ppginf/book 12pt, textwidth 16cm):
% largura total ~15,6 cm, altura ~5 cm (fração de página, não página inteira).
% Legível em P&B: etapas = cinza claro; destino = cinza médio; LCE tracejada.
% ============================================================================
\begin{tikzpicture}[
    x=1cm, y=1cm,
    etapa/.style={draw, rounded corners=2pt, align=center, inner sep=2.5pt,
                  font=\scriptsize, fill=gray!8, minimum height=17mm},
    destino/.style={etapa, fill=gray!16},
    seta/.style={-{Stealth[length=2.2mm]}, semithick},
    rot/.style={font=\tiny, align=center, inner sep=1pt}
]
    % ------------------ a linha do processo, da esquerda para a direita ----
    \node[etapa, text width=20mm] (prob) at (0,0)
        {\textbf{problema}\\(Cap.~\ref{ch:intro}):\\reduzir o custo de\\rotulagem sem perda\\de desempenho};
    \node[etapa, text width=21mm] (terr) at (2.68,0)
        {\textbf{terreno}\\\textbf{comum}:\\desenho e registro;\\dados e auditoria;\\classificadores;\\métricas};
    \node[etapa, text width=20mm] (p1) at (5.36,0)
        {\textbf{pilar 1}:\\quanto importa a\\composição do\\conjunto inicial $L_0$};
    \node[etapa, text width=20mm] (p2) at (8.02,0)
        {\textbf{pilar 2}:\\construir $L_0$\\sem rótulos\\(DRI-SL)};
    \node[etapa, text width=20mm] (p3) at (10.66,0)
        {\textbf{pilar 3}:\\quem rotula:\\LLM como oráculo\\(acerto, custo, erro)};
    \node[destino, text width=21mm] (p4) at (13.42,0)
        {\textbf{pilar 4}:\\integração no\\FALCO e avaliação\\sob o mesmo\\orçamento};

    % ------------------ setas com a lógica de precedência ------------------
    \draw[seta] (prob) -- (terr);
    \draw[seta] (terr) -- (p1);
    \draw[seta] (p1) -- (p2);
    \draw[seta] (p2) -- (p3);
    \draw[seta] (p3) -- (p4);
    \node[rot, anchor=south] at (4.02,1.3) {o processo começa\\sem rótulo algum};
    \node[rot, anchor=south] at (6.69,1.3)  {medida a importância,\\constrói-se o começo};
    \node[rot, anchor=south] at (9.34,1.3)  {definido o começo:\\quem rotula?};
    \node[rot, anchor=south] at (12.04,1.3)  {resolvidos início\\e oráculo};

    % ------------------ LCE transversal + fecho, abaixo da linha -----------
    \draw[dashed, rounded corners=2pt]
        (4.2,-1.55) rectangle (14.55,-2.35);
    \node[font=\scriptsize, align=center] at (9.37,-1.95)
        {\textbf{transversal}: a LCE resume as curvas de aprendizado dos quatro pilares};
    \node[font=\tiny, align=center, anchor=north] at (7.25,-2.6)
        {o capítulo encerra com a síntese do desenho experimental};
\end{tikzpicture}

% SEM números medidos, SEM códigos internos, SEM caminhos, SEM travessões.
% Uso standalone (pré-visualização):
%   \documentclass[tikz,border=6pt,12pt]{standalone}
%   \usetikzlibrary{arrows.meta, positioning}
%   \begin{document}% ============================================================================
% ESQUEMA: a sequência da metodologia como fluxo de decisões (HORIZONTAL v2).
% Encomenda do autor (tarefa 0138; reorientada na auditoria de 2026-08-24:
% "prefiro que fique na horizontal e menor, está ocupando uma página inteira").
% Conta a história da abertura do Cap. 3 na mesma ordem: problema -> terreno
% comum -> pilares 1-4 (setas com a lógica de precedência) -> LCE transversal
% -> fecho (validade e reprodutibilidade).
% Uso: abertura do Cap. 3, dentro de figure:
%   % ============================================================================
% ESQUEMA: a sequência da metodologia como fluxo de decisões (HORIZONTAL v2).
% Encomenda do autor (tarefa 0138; reorientada na auditoria de 2026-08-24:
% "prefiro que fique na horizontal e menor, está ocupando uma página inteira").
% Conta a história da abertura do Cap. 3 na mesma ordem: problema -> terreno
% comum -> pilares 1-4 (setas com a lógica de precedência) -> LCE transversal
% -> fecho (validade e reprodutibilidade).
% Uso: abertura do Cap. 3, dentro de figure:
%   \input{3-metodo/esquemas-propostos/esq-sequencia-metodologia}
% SEM números medidos, SEM códigos internos, SEM caminhos, SEM travessões.
% Uso standalone (pré-visualização):
%   \documentclass[tikz,border=6pt,12pt]{standalone}
%   \usetikzlibrary{arrows.meta, positioning}
%   \begin{document}\input{esq-sequencia-metodologia.tex}\end{document}
%   (no modo standalone, troque a \ref{ch:intro} por texto fixo)
% GEOMETRIA CALIBRADA PARA CORPO 12 (ppginf/book 12pt, textwidth 16cm):
% largura total ~15,6 cm, altura ~5 cm (fração de página, não página inteira).
% Legível em P&B: etapas = cinza claro; destino = cinza médio; LCE tracejada.
% ============================================================================
\begin{tikzpicture}[
    x=1cm, y=1cm,
    etapa/.style={draw, rounded corners=2pt, align=center, inner sep=2.5pt,
                  font=\scriptsize, fill=gray!8, minimum height=17mm},
    destino/.style={etapa, fill=gray!16},
    seta/.style={-{Stealth[length=2.2mm]}, semithick},
    rot/.style={font=\tiny, align=center, inner sep=1pt}
]
    % ------------------ a linha do processo, da esquerda para a direita ----
    \node[etapa, text width=20mm] (prob) at (0,0)
        {\textbf{problema}\\(Cap.~\ref{ch:intro}):\\reduzir o custo de\\rotulagem sem perda\\de desempenho};
    \node[etapa, text width=21mm] (terr) at (2.68,0)
        {\textbf{terreno}\\\textbf{comum}:\\desenho e registro;\\dados e auditoria;\\classificadores;\\métricas};
    \node[etapa, text width=20mm] (p1) at (5.36,0)
        {\textbf{pilar 1}:\\quanto importa a\\composição do\\conjunto inicial $L_0$};
    \node[etapa, text width=20mm] (p2) at (8.02,0)
        {\textbf{pilar 2}:\\construir $L_0$\\sem rótulos\\(DRI-SL)};
    \node[etapa, text width=20mm] (p3) at (10.66,0)
        {\textbf{pilar 3}:\\quem rotula:\\LLM como oráculo\\(acerto, custo, erro)};
    \node[destino, text width=21mm] (p4) at (13.42,0)
        {\textbf{pilar 4}:\\integração no\\FALCO e avaliação\\sob o mesmo\\orçamento};

    % ------------------ setas com a lógica de precedência ------------------
    \draw[seta] (prob) -- (terr);
    \draw[seta] (terr) -- (p1);
    \draw[seta] (p1) -- (p2);
    \draw[seta] (p2) -- (p3);
    \draw[seta] (p3) -- (p4);
    \node[rot, anchor=south] at (4.02,1.3) {o processo começa\\sem rótulo algum};
    \node[rot, anchor=south] at (6.69,1.3)  {medida a importância,\\constrói-se o começo};
    \node[rot, anchor=south] at (9.34,1.3)  {definido o começo:\\quem rotula?};
    \node[rot, anchor=south] at (12.04,1.3)  {resolvidos início\\e oráculo};

    % ------------------ LCE transversal + fecho, abaixo da linha -----------
    \draw[dashed, rounded corners=2pt]
        (4.2,-1.55) rectangle (14.55,-2.35);
    \node[font=\scriptsize, align=center] at (9.37,-1.95)
        {\textbf{transversal}: a LCE resume as curvas de aprendizado dos quatro pilares};
    \node[font=\tiny, align=center, anchor=north] at (7.25,-2.6)
        {o capítulo encerra com a síntese do desenho experimental};
\end{tikzpicture}

% SEM números medidos, SEM códigos internos, SEM caminhos, SEM travessões.
% Uso standalone (pré-visualização):
%   \documentclass[tikz,border=6pt,12pt]{standalone}
%   \usetikzlibrary{arrows.meta, positioning}
%   \begin{document}% ============================================================================
% ESQUEMA: a sequência da metodologia como fluxo de decisões (HORIZONTAL v2).
% Encomenda do autor (tarefa 0138; reorientada na auditoria de 2026-08-24:
% "prefiro que fique na horizontal e menor, está ocupando uma página inteira").
% Conta a história da abertura do Cap. 3 na mesma ordem: problema -> terreno
% comum -> pilares 1-4 (setas com a lógica de precedência) -> LCE transversal
% -> fecho (validade e reprodutibilidade).
% Uso: abertura do Cap. 3, dentro de figure:
%   \input{3-metodo/esquemas-propostos/esq-sequencia-metodologia}
% SEM números medidos, SEM códigos internos, SEM caminhos, SEM travessões.
% Uso standalone (pré-visualização):
%   \documentclass[tikz,border=6pt,12pt]{standalone}
%   \usetikzlibrary{arrows.meta, positioning}
%   \begin{document}\input{esq-sequencia-metodologia.tex}\end{document}
%   (no modo standalone, troque a \ref{ch:intro} por texto fixo)
% GEOMETRIA CALIBRADA PARA CORPO 12 (ppginf/book 12pt, textwidth 16cm):
% largura total ~15,6 cm, altura ~5 cm (fração de página, não página inteira).
% Legível em P&B: etapas = cinza claro; destino = cinza médio; LCE tracejada.
% ============================================================================
\begin{tikzpicture}[
    x=1cm, y=1cm,
    etapa/.style={draw, rounded corners=2pt, align=center, inner sep=2.5pt,
                  font=\scriptsize, fill=gray!8, minimum height=17mm},
    destino/.style={etapa, fill=gray!16},
    seta/.style={-{Stealth[length=2.2mm]}, semithick},
    rot/.style={font=\tiny, align=center, inner sep=1pt}
]
    % ------------------ a linha do processo, da esquerda para a direita ----
    \node[etapa, text width=20mm] (prob) at (0,0)
        {\textbf{problema}\\(Cap.~\ref{ch:intro}):\\reduzir o custo de\\rotulagem sem perda\\de desempenho};
    \node[etapa, text width=21mm] (terr) at (2.68,0)
        {\textbf{terreno}\\\textbf{comum}:\\desenho e registro;\\dados e auditoria;\\classificadores;\\métricas};
    \node[etapa, text width=20mm] (p1) at (5.36,0)
        {\textbf{pilar 1}:\\quanto importa a\\composição do\\conjunto inicial $L_0$};
    \node[etapa, text width=20mm] (p2) at (8.02,0)
        {\textbf{pilar 2}:\\construir $L_0$\\sem rótulos\\(DRI-SL)};
    \node[etapa, text width=20mm] (p3) at (10.66,0)
        {\textbf{pilar 3}:\\quem rotula:\\LLM como oráculo\\(acerto, custo, erro)};
    \node[destino, text width=21mm] (p4) at (13.42,0)
        {\textbf{pilar 4}:\\integração no\\FALCO e avaliação\\sob o mesmo\\orçamento};

    % ------------------ setas com a lógica de precedência ------------------
    \draw[seta] (prob) -- (terr);
    \draw[seta] (terr) -- (p1);
    \draw[seta] (p1) -- (p2);
    \draw[seta] (p2) -- (p3);
    \draw[seta] (p3) -- (p4);
    \node[rot, anchor=south] at (4.02,1.3) {o processo começa\\sem rótulo algum};
    \node[rot, anchor=south] at (6.69,1.3)  {medida a importância,\\constrói-se o começo};
    \node[rot, anchor=south] at (9.34,1.3)  {definido o começo:\\quem rotula?};
    \node[rot, anchor=south] at (12.04,1.3)  {resolvidos início\\e oráculo};

    % ------------------ LCE transversal + fecho, abaixo da linha -----------
    \draw[dashed, rounded corners=2pt]
        (4.2,-1.55) rectangle (14.55,-2.35);
    \node[font=\scriptsize, align=center] at (9.37,-1.95)
        {\textbf{transversal}: a LCE resume as curvas de aprendizado dos quatro pilares};
    \node[font=\tiny, align=center, anchor=north] at (7.25,-2.6)
        {o capítulo encerra com a síntese do desenho experimental};
\end{tikzpicture}
\end{document}
%   (no modo standalone, troque a \ref{ch:intro} por texto fixo)
% GEOMETRIA CALIBRADA PARA CORPO 12 (ppginf/book 12pt, textwidth 16cm):
% largura total ~15,6 cm, altura ~5 cm (fração de página, não página inteira).
% Legível em P&B: etapas = cinza claro; destino = cinza médio; LCE tracejada.
% ============================================================================
\begin{tikzpicture}[
    x=1cm, y=1cm,
    etapa/.style={draw, rounded corners=2pt, align=center, inner sep=2.5pt,
                  font=\scriptsize, fill=gray!8, minimum height=17mm},
    destino/.style={etapa, fill=gray!16},
    seta/.style={-{Stealth[length=2.2mm]}, semithick},
    rot/.style={font=\tiny, align=center, inner sep=1pt}
]
    % ------------------ a linha do processo, da esquerda para a direita ----
    \node[etapa, text width=20mm] (prob) at (0,0)
        {\textbf{problema}\\(Cap.~\ref{ch:intro}):\\reduzir o custo de\\rotulagem sem perda\\de desempenho};
    \node[etapa, text width=21mm] (terr) at (2.68,0)
        {\textbf{terreno}\\\textbf{comum}:\\desenho e registro;\\dados e auditoria;\\classificadores;\\métricas};
    \node[etapa, text width=20mm] (p1) at (5.36,0)
        {\textbf{pilar 1}:\\quanto importa a\\composição do\\conjunto inicial $L_0$};
    \node[etapa, text width=20mm] (p2) at (8.02,0)
        {\textbf{pilar 2}:\\construir $L_0$\\sem rótulos\\(DRI-SL)};
    \node[etapa, text width=20mm] (p3) at (10.66,0)
        {\textbf{pilar 3}:\\quem rotula:\\LLM como oráculo\\(acerto, custo, erro)};
    \node[destino, text width=21mm] (p4) at (13.42,0)
        {\textbf{pilar 4}:\\integração no\\FALCO e avaliação\\sob o mesmo\\orçamento};

    % ------------------ setas com a lógica de precedência ------------------
    \draw[seta] (prob) -- (terr);
    \draw[seta] (terr) -- (p1);
    \draw[seta] (p1) -- (p2);
    \draw[seta] (p2) -- (p3);
    \draw[seta] (p3) -- (p4);
    \node[rot, anchor=south] at (4.02,1.3) {o processo começa\\sem rótulo algum};
    \node[rot, anchor=south] at (6.69,1.3)  {medida a importância,\\constrói-se o começo};
    \node[rot, anchor=south] at (9.34,1.3)  {definido o começo:\\quem rotula?};
    \node[rot, anchor=south] at (12.04,1.3)  {resolvidos início\\e oráculo};

    % ------------------ LCE transversal + fecho, abaixo da linha -----------
    \draw[dashed, rounded corners=2pt]
        (4.2,-1.55) rectangle (14.55,-2.35);
    \node[font=\scriptsize, align=center] at (9.37,-1.95)
        {\textbf{transversal}: a LCE resume as curvas de aprendizado dos quatro pilares};
    \node[font=\tiny, align=center, anchor=north] at (7.25,-2.6)
        {o capítulo encerra com a síntese do desenho experimental};
\end{tikzpicture}
\end{document}
%   (no modo standalone, troque a \ref{ch:intro} por texto fixo)
% GEOMETRIA CALIBRADA PARA CORPO 12 (ppginf/book 12pt, textwidth 16cm):
% largura total ~15,6 cm, altura ~5 cm (fração de página, não página inteira).
% Legível em P&B: etapas = cinza claro; destino = cinza médio; LCE tracejada.
% ============================================================================
\begin{tikzpicture}[
    x=1cm, y=1cm,
    etapa/.style={draw, rounded corners=2pt, align=center, inner sep=2.5pt,
                  font=\scriptsize, fill=gray!8, minimum height=17mm},
    destino/.style={etapa, fill=gray!16},
    seta/.style={-{Stealth[length=2.2mm]}, semithick},
    rot/.style={font=\tiny, align=center, inner sep=1pt}
]
    % ------------------ a linha do processo, da esquerda para a direita ----
    \node[etapa, text width=20mm] (prob) at (0,0)
        {\textbf{problema}\\(Cap.~\ref{ch:intro}):\\reduzir o custo de\\rotulagem sem perda\\de desempenho};
    \node[etapa, text width=21mm] (terr) at (2.68,0)
        {\textbf{terreno}\\\textbf{comum}:\\desenho e registro;\\dados e auditoria;\\classificadores;\\métricas};
    \node[etapa, text width=20mm] (p1) at (5.36,0)
        {\textbf{pilar 1}:\\quanto importa a\\composição do\\conjunto inicial $L_0$};
    \node[etapa, text width=20mm] (p2) at (8.02,0)
        {\textbf{pilar 2}:\\construir $L_0$\\sem rótulos\\(DRI-SL)};
    \node[etapa, text width=20mm] (p3) at (10.66,0)
        {\textbf{pilar 3}:\\quem rotula:\\LLM como oráculo\\(acerto, custo, erro)};
    \node[destino, text width=21mm] (p4) at (13.42,0)
        {\textbf{pilar 4}:\\integração no\\FALCO e avaliação\\sob o mesmo\\orçamento};

    % ------------------ setas com a lógica de precedência ------------------
    \draw[seta] (prob) -- (terr);
    \draw[seta] (terr) -- (p1);
    \draw[seta] (p1) -- (p2);
    \draw[seta] (p2) -- (p3);
    \draw[seta] (p3) -- (p4);
    \node[rot, anchor=south] at (4.02,1.3) {o processo começa\\sem rótulo algum};
    \node[rot, anchor=south] at (6.69,1.3)  {medida a importância,\\constrói-se o começo};
    \node[rot, anchor=south] at (9.34,1.3)  {definido o começo:\\quem rotula?};
    \node[rot, anchor=south] at (12.04,1.3)  {resolvidos início\\e oráculo};

    % ------------------ LCE transversal + fecho, abaixo da linha -----------
    \draw[dashed, rounded corners=2pt]
        (4.2,-1.55) rectangle (14.55,-2.35);
    \node[font=\scriptsize, align=center] at (9.37,-1.95)
        {\textbf{transversal}: a LCE resume as curvas de aprendizado dos quatro pilares};
    \node[font=\tiny, align=center, anchor=north] at (7.25,-2.6)
        {o capítulo encerra com a síntese do desenho experimental};
\end{tikzpicture}
\end{document}
%   (no modo standalone, troque a \ref{ch:intro} por texto fixo)
% GEOMETRIA CALIBRADA PARA CORPO 12 (ppginf/book 12pt, textwidth 16cm):
% largura total ~15,6 cm, altura ~5 cm (fração de página, não página inteira).
% Legível em P&B: etapas = cinza claro; destino = cinza médio; LCE tracejada.
% ============================================================================
\begin{tikzpicture}[
    x=1cm, y=1cm,
    etapa/.style={draw, rounded corners=2pt, align=center, inner sep=2.5pt,
                  font=\scriptsize, fill=gray!8, minimum height=17mm},
    destino/.style={etapa, fill=gray!16},
    seta/.style={-{Stealth[length=2.2mm]}, semithick},
    rot/.style={font=\tiny, align=center, inner sep=1pt}
]
    % ------------------ a linha do processo, da esquerda para a direita ----
    \node[etapa, text width=20mm] (prob) at (0,0)
        {\textbf{problema}\\(Cap.~\ref{ch:intro}):\\reduzir o custo de\\rotulagem sem perda\\de desempenho};
    \node[etapa, text width=21mm] (terr) at (2.68,0)
        {\textbf{terreno}\\\textbf{comum}:\\desenho e registro;\\dados e auditoria;\\classificadores;\\métricas};
    \node[etapa, text width=20mm] (p1) at (5.36,0)
        {\textbf{pilar 1}:\\quanto importa a\\composição do\\conjunto inicial $L_0$};
    \node[etapa, text width=20mm] (p2) at (8.02,0)
        {\textbf{pilar 2}:\\construir $L_0$\\sem rótulos\\(DRI-SL)};
    \node[etapa, text width=20mm] (p3) at (10.66,0)
        {\textbf{pilar 3}:\\quem rotula:\\LLM como oráculo\\(acerto, custo, erro)};
    \node[destino, text width=21mm] (p4) at (13.42,0)
        {\textbf{pilar 4}:\\integração no\\FALCO e avaliação\\sob o mesmo\\orçamento};

    % ------------------ setas com a lógica de precedência ------------------
    \draw[seta] (prob) -- (terr);
    \draw[seta] (terr) -- (p1);
    \draw[seta] (p1) -- (p2);
    \draw[seta] (p2) -- (p3);
    \draw[seta] (p3) -- (p4);
    \node[rot, anchor=south] at (4.02,1.3) {o processo começa\\sem rótulo algum};
    \node[rot, anchor=south] at (6.69,1.3)  {medida a importância,\\constrói-se o começo};
    \node[rot, anchor=south] at (9.34,1.3)  {definido o começo:\\quem rotula?};
    \node[rot, anchor=south] at (12.04,1.3)  {resolvidos início\\e oráculo};

    % ------------------ LCE transversal + fecho, abaixo da linha -----------
    \draw[dashed, rounded corners=2pt]
        (4.2,-1.55) rectangle (14.55,-2.35);
    \node[font=\scriptsize, align=center] at (9.37,-1.95)
        {\textbf{transversal}: a LCE resume as curvas de aprendizado dos quatro pilares};
    \node[font=\tiny, align=center, anchor=north] at (7.25,-2.6)
        {o capítulo encerra com a síntese do desenho experimental};
\end{tikzpicture}
